%Regarding my main concern---spillovers---the current version is considerably more transparent, and Section 4.2, together with the Online Appendix, attempts to address this threat to identification. I would, however, encourage the authors to reorganize the presentation of this evidence. My suggestion would be to lead with the exposure-based specification, in which exposure to treated neighbors is measured continuously. In my view, this is the most informative of the three exercises because it exploits variation in the intensity of potential spillovers rather than relying on a binary indicator for having a treated neighbor. This specification should, of course, be presented with the appropriate caveat: exposure to treated neighbors is not randomly assigned, so the estimates do not fully address the possibility that municipalities with many early-adopting neighbors differ systematically in other respects.


Thanks. We have reorganized the presentation of the evidence on spillovers, as suggested. Section 4.2 now leads with the exposure-based specification following \citet{crepon2013}, which measures exposure to treated neighbors continuously (as the share and the number of treated neighbors), and only then presents the binary neighbor-treatment specification. As you note, the exposure-based specification is the more informative of the two, since it exploits variation in the intensity of potential spillovers rather than relying on a binary indicator for having a treated neighbor.

When commenting on the results of the exposure-based specification, we have introduced the recommended caveat, which now reads as follows in the paper: ``Because exposure to treated neighbors is not randomly assigned, however, this evidence does not fully rule out the possibility that municipalities with many early-adopting neighbors differ systematically from those with fewer such neighbors in other respects that are also correlated with our outcomes.''

Separately, and to keep the spillovers discussion focused, we have also removed the randomization-inference exercise from this section, since (as discussed in our response to your next comment) it is better understood as a general robustness check on our main estimates than as a targeted test for spillovers. We now report it separately when discussing our main results, with the full exercise moved to its own subsection in the Online Appendix (Section C.8).